% !TeX root = ../../Book3.tex
% 纲要标题：### 2.3 模的支集
% 纲要位置：工作记录/计划纲要.md，第 1955 行。
% 纲要细目（写作范围）：
% * 支集、零化子与有限生成性；循环模支集等于理想闭集
% * 局部化与 Hom 的交换条件；有限生成不能代替有限展示的循环模反例
% * 支集与短正合列

\section{模的支集}
\label{b3:sec:0203}

局部化把一个模分散到素谱的各点；支集记录它在哪些点仍然非零。有限生成性使这些局部信息受到一个理想的统一控制。处理同态时，还要进一步区分“有限个生成元”与“有限个生成关系”，因为清分母必须一次覆盖全部有关数据。

\subsection{支集、零化子与有限生成性}
\label{b3:subsec:0203:01}

\begin{definition}\label{b3:def:support-annihilator}
设 \(R\) 为交换含幺环，\(M\) 为 \(R\) 模。\(M\) 的\term[zhiji]{支集}[support]与\term[linghuazi]{零化子}[annihilator]分别为
\[
\begin{aligned}
\operatorname{Supp}_R M
&=\{\,\mathfrak p\in\operatorname{Spec}R\mid M_{\mathfrak p}\ne0\,\},\\
\operatorname{Ann}_R M
&=\{\,a\in R\mid am=0\text{ 对所有 }m\in M\,\}.
\end{aligned}
\]
零化子也称湮灭子。\termidx[yinmiezi]{湮灭子}对单个元素，简记
\(\operatorname{Ann}_R(m)=\{a\in R:am=0\}\)；它等于循环子模 \(Rm\) 的零化子。
\end{definition}
\symidx{\(\operatorname{Supp}_R M\)：模的支集}
\symidx{\(\operatorname{Ann}_R M\)：模的零化子}

\begin{theorem}\label{b3:thm:support-finite}
设 \(R\) 为交换含幺环，\(M\) 为 \(R\) 模。则
\[
\operatorname{Supp}_R M\subseteq V(\operatorname{Ann}_R M).
\]
若 \(M\) 有限生成，则等号成立。更具体地，对每个 \(R\) 的素理想 \(\mathfrak p\)，有限生成模 \(M\) 满足
\[
M_{\mathfrak p}=0
\quad\Longleftrightarrow\quad
sM=0\text{ 对某个 }s\notin\mathfrak p\text{ 成立}.
\]
\end{theorem}
\begin{proof}
若 \(a\in\operatorname{Ann}_R M\setminus\mathfrak p\)，在 \(R_{\mathfrak p}\) 中 \(a\) 可逆，又零化 \(M_{\mathfrak p}\)，故该模为零。这证明一般包含关系。

设 \(M=Rm_1+\cdots+Rm_r\) 且 \(M_{\mathfrak p}=0\)。由\eqref{b3:eq:localized-element-zero}，对每个 \(i\) 有 \(s_i\notin\mathfrak p\) 使 \(s_im_i=0\)。取 \(s=s_1\cdots s_r\)，则 \(s\notin\mathfrak p\) 且 \(sM=0\)。反向由 \(s\) 局部可逆得到。于是 \(M_{\mathfrak p}=0\) 等价于 \(\operatorname{Ann}_R M\not\subseteq\mathfrak p\)，即所述等式。对零模取空生成集及 \(s=1\)，论证仍成立。
\end{proof}

因而有限生成模的支集是闭集。特别地，对 \(R\) 的任意理想 \(I\)，循环模 \(R/I\) 有限生成，且 \(\operatorname{Ann}_R(R/I)=I\)：零化它的元素必须零化 \(1+I\)，反向则由商模的运算得到。因此
\[
\operatorname{Supp}_R(R/I)=V(I).
\]
例如，对域 \(k\)，有
\[
\operatorname{Supp}_{k[x,y]} k[x,y]/(xy)=V(xy)=V(x)\cup V(y).
\]
模掉 \(xy\) 得到的模同时分布在两条坐标轴上；若进一步在 \(x\) 处局部化，则只留下 \(y=0\) 且 \(x\ne0\) 的部分。

\ProseParagraphBreak
有限生成条件不能省去。整数模 \(\mathbb Q/\mathbb Z\) 的零化子为零：对非零整数 \(n\)，选不能整除 \(n\) 的正整数 \(d\)，则 \(n(1/d+\mathbb Z)\ne0\)。但在 \((0)\) 处局部化后，每个元素都被某个非零整数零化，故整个模变成零。另一方面，对任一素数 \(p\)，元素 \(1/p+\mathbb Z\) 在 \((p)\) 处局部化后仍不为零：若它变成零，就有某个 \(s\in\mathbb Z\setminus(p)\) 使 \(s/p\in\mathbb Z\)，与 \(p\nmid s\) 矛盾。由于 \(\mathbb Z\) 的素理想只有 \((0)\) 及各个 \((p)\)，这就完整算出
\[
\operatorname{Supp}_{\mathbb Z}(\mathbb Q/\mathbb Z)
=\{\,(p)\mid p\text{ 为素数}\,\},\qquad
V\bigl(\operatorname{Ann}_{\mathbb Z}(\mathbb Q/\mathbb Z)\bigr)
=\operatorname{Spec}\mathbb Z.
\]
非零整数不可能被所有素数整除，故 \(\bigcap_{p\text{ 为素数}}(p)=(0)\)。若闭集 \(V(I)\) 包含上述支集，则 \(I\) 包含于这个交，只能为零。因此支集在 \(\operatorname{Spec}\mathbb Z\) 中稠密；它又缺少 \((0)\)，因而不是闭集。此例中每个元素各有一个零化它的非零乘子，却没有一个非零整数同时零化全部元素。

\subsection{局部化与 Hom 的交换条件}
\label{b3:subsec:0203:02}

设 \(R\) 为交换含幺环，\(S\subseteq R\) 为含一的乘法集，\(M,N\) 为 \(R\) 模。局部化后的同态可以对分式取值，因而有自然映射
\begin{equation}\label{b3:eq:hom-localization-map}
S^{-1}\operatorname{Hom}_R(M,N)
\longrightarrow
\operatorname{Hom}_{S^{-1}R}(S^{-1}M,S^{-1}N),
\qquad
\frac fs\longmapsto
\left(\frac mt\longmapsto\frac{f(m)}{st}\right).
\end{equation}
若两侧的分式代表元改变，交叉差仍被某个分母零化；因此该公式良定义。也可先在 \(M\) 上取 \(m\mapsto f(m)/s\)，再用\cref{b3:prop:module-fraction-operations}的延拓泛性质。

\begin{theorem}\label{b3:thm:hom-localization}
设 \(R\) 为交换含幺环，\(S\subseteq R\) 为含一的乘法集，\(M,N\) 为 \(R\) 模。考虑自然映射
\[
S^{-1}\operatorname{Hom}_R(M,N)\longrightarrow
\operatorname{Hom}_{S^{-1}R}(S^{-1}M,S^{-1}N),\qquad
f/s\longmapsto\bigl(m/t\longmapsto f(m)/(st)\bigr).
\]
若 \(M\) 有限生成，则此映射为单射。若 \(M\) 有限展示，即存在非负整数 \(a,b\) 及 \(R\) 模正合列
\[
R^a\longrightarrow R^b\longrightarrow M\longrightarrow0,
\]
则该映射为同构。两种结论都不要求 \(N\) 有限生成。
\end{theorem}
\begin{proof}
先只假设 \(M\) 有限生成，取生成元 \(m_1,\ldots,m_r\)。单射性要求：一个同态若在局部化后消失，就能找到一个分母零化这个同态的全部值。若 \(f/s\) 在右侧为零，则每个 \(f(m_i)/s=0\)，故可取 \(u_i\in S\) 使 \(u_i\cdot f(m_i)=0\)。因为生成元只有有限个，乘积 \(u=\prod_{i=1}^r u_i\) 零化 \(f\) 在全部生成元上的值，因而 \(uf=0\)，即 \(f/s=0\)。这证明单射性；当 \(M=0\) 时可取空生成集及 \(u=1\)。

再设 \(M\) 有上述有限展示。要把局部化后的同态写成 \(f/s\)，仅给生成元的像通分还不够；通分后的像必须在 \(N\) 中满足全部关系，才能定义 \(f:M\to N\)。记 \(e_1,\ldots,e_b\) 在 \(M\) 中的像为 \(m_i\)。展示的第一个映射的像由有限个关系
\(\sum_i a_{ij}m_i=0\)、\(1\le j\le a\) 生成。给定右侧的同态 \(g\)，把有限个 \(g(m_i/1)\) 通分为 \(n_i/s\)。每个关系给出
\[
\frac{\sum_i a_{ij}n_i}{s}=0.
\]
故可选一个共同的 \(u\in S\)，使全部 \(\sum_i a_{ij}u n_i=0\) 在 \(N\) 中成立。于是 \(m_i\mapsto u n_i\) 满足全部关系，诱导 \(f:M\rightarrow N\)。此时 \(f/(us)\) 在\eqref{b3:eq:hom-localization-map}下的像在 \(m_i/1\) 上等于 \(g\)，因而就是 \(g\)。这证明满射性。
\end{proof}

有限生成不能代替有限展示来保证满射性。下面的循环模只有一个生成元，却无法同时清除所有关系中的分母。

\begin{proposition}\label{b3:prop:hom-localization-fg-counterexample}
设 \(k\) 为域，\(R=\prod_{n\ge1}k\)，\(I=\bigoplus_{n\ge1}k\subseteq R\) 为只有有限多个非零坐标的元素组成的理想。对有限集 \(F\subseteq\mathbb Z_{\ge1}\)，令 \(e_F\) 在 \(F\) 的坐标上等于 \(1\)，其余坐标上等于 \(0\)，并取 \(S=\{1-e_F:F\text{ 有限}\}\)。则 \(S\) 是含一而不含零的乘法集，\(M=R/I\) 是有限生成但非有限展示的 \(R\) 模，且自然映射
\[
S^{-1}\operatorname{Hom}_R(M,R)\longrightarrow
\operatorname{Hom}_{S^{-1}R}(S^{-1}M,S^{-1}R)
\]
的源为零，目标非零。还有自然同构 \(S^{-1}R\cong R/I\)，把 \(a/s\) 送到 \(a+I\)。
\end{proposition}
\begin{proof}
因为 \((1-e_F)(1-e_G)=1-e_{F\cup G}\)，\(S\) 对乘法封闭；取空集得到 \(1\in S\)，而有限集不可能包含所有正整数，所以 \(0\notin S\)。\(M\) 由 \(1+I\) 生成。一个同态 \(M\to R\) 在该生成元上的值 \(a\) 必须满足 \(Ia=0\)。用每个单坐标幂等元 \(e_{\{n\}}\in I\) 相乘得到 \(a_n=0\)，故 \(a=0\)，即
\[
\operatorname{Hom}_R(M,R)=\operatorname{Ann}_R I=0.
\]
另一方面，每个 \(x\in I\) 只有有限多个非零坐标；若这些坐标包含于 \(F\)，则 \((1-e_F)x=0\)。由分式零判据，\(S^{-1}I=0\)，局部化正合性于是给出 \(S^{-1}M\cong S^{-1}R\)。环 \(S^{-1}R\) 非零，因为 \(1/1=0\) 将要求 \(S\) 含有零元。因此目标包含非零恒等同态，而源为零，映射不满射。由\cref{b3:thm:hom-localization}，\(M\) 不可能有限展示。

对每个 \(s=1-e_F\in S\)，其在 \(R/I\) 中的像等于 \(1\)，故商映射延拓为 \(S^{-1}R\to R/I\)、\(a/s\mapsto a+I\)。其逆为 \(a+I\mapsto a/1\)：\(S^{-1}I=0\) 保证该映射良定义，且 \(e_F/1=0\) 使 \(s/1=1\)，从而 \(a/s=a/1\)。
\end{proof}

若连有限生成也不假设，取 \(R=\mathbb Z\)、\(S=\mathbb Z\setminus\{0\}\)、\(M=\mathbb Q\)、\(N=\mathbb Z\)。对任意 \(\mathbb Z\) 模同态 \(f:\mathbb Q\rightarrow\mathbb Z\)，有
\[
f(q)=n\cdot f(q/n)\qquad(q\in\mathbb Q,\ n\ge1).
\]
所以整数 \(f(q)\) 被每个正整数整除，只能为零。因而 \(\operatorname{Hom}_{\mathbb Z}(\mathbb Q,\mathbb Z)=0\)，而局部化后两模都成为 \(\mathbb Q\)，右侧为 \(\operatorname{Hom}_{\mathbb Q}(\mathbb Q,\mathbb Q)=\mathbb Q\)。因此一般的 Hom 不能无条件移过局部化。

\subsection{短正合列中的支集}
\label{b3:subsec:0203:03}

\begin{proposition}\label{b3:prop:support-exact}
设 \(R\) 为交换含幺环，\(L,M,N\) 为 \(R\) 模。若 \(0\rightarrow L\rightarrow M\rightarrow N\rightarrow0\) 正合，则
\[
\operatorname{Supp}_R M
=\operatorname{Supp}_R L\cup\operatorname{Supp}_R N.
\]
该结论对任意 \(R\) 模成立。此外，任意一族 \(R\) 模 \((M_\lambda)_\lambda\) 满足
\[
\operatorname{Supp}_R\left(\bigoplus_{\lambda}M_\lambda\right)
=\bigcup_{\lambda}\operatorname{Supp}_R M_\lambda.
\]
\end{proposition}
\begin{proof}
由\cref{b3:thm:localization-exact}，在每个 \(\mathfrak p\) 处仍有短正合列。中间项为零当且仅当两端同时为零，取补集即得第一式。

局部化与直和相容：把分式 \((m_\lambda)_\lambda/s\) 送到 \((m_\lambda/s)_\lambda\)。其像仍只有有限多个非零分量；反过来，对这些非零分量的分母取乘积，就能通分并得到逆映射。这里的有限性来自直和的定义，与素谱上模的支集是不同的概念。直和为零当且仅当各分量为零，因而得到第二式。
\end{proof}

第一式说明支集忽略了短正合列是否分裂。例如，设 \(k\) 为域，在 \(R=k[x]\) 上有
\[
0\longrightarrow(x)/(x^2)\longrightarrow R/(x^2)
\longrightarrow R/(x)\longrightarrow0
\]
的三个非零模支集都为 \(\bigl\{(x)\bigr\}\)，但中间项中的 \(x\) 作用非零而平方为零，两端的 \(x\) 作用均为零。这种厚度信息并未由支集单独记录，后面还要用长度、关联素理想及准素分解作更细的分析。
